\section{Архитектура системы на основе PINN}

\subsection{Общая структура}

Компоненты системы:

\begin{enumerate}
    \item Модуль предобработки данных
    \item PINN (нейронная сеть)
    \item Модуль обучения
    \item Модуль постобработки
    \item Модуль управления
\end{enumerate}

\subsection{Архитектура PINN}

\subsubsection{Структура нейронной сети}

PINN аппроксимирует решение:
\begin{equation}
\hat{u}(x, y, z, t; \theta), \quad \hat{T}(x, y, z, t; \theta)
\end{equation}

Архитектура:
\begin{itemize}
    \item Вход: $(x, y, z, t)$
    \item Скрытые слои: полносвязные с активацией $\sigma$
    \item Выход: $(u_x, u_y, u_z, T)$
\end{itemize}

\begin{equation}
\begin{aligned}
h_0 &= [x, y, z, t]^T \\
h_l &= \sigma(W_l h_{l-1} + b_l), \quad l = 1, \ldots, L-1 \\
\hat{u}, \hat{T} &= W_L h_{L-1} + b_L
\end{aligned}
\end{equation}

\subsubsection{Функция потерь}

\begin{equation}
\mathcal{L}_{total} = \mathcal{L}_{PDE} + \lambda_{BC}\mathcal{L}_{BC} + \lambda_{IC}\mathcal{L}_{IC} + \lambda_{data}\mathcal{L}_{data}
\end{equation}

Остатки PDE:

\begin{align}
\mathcal{R}_1 &= \rho \frac{\partial^2 \hat{u}_j}{\partial t^2} - \frac{\partial \hat{\sigma}_{ij}}{\partial x_i} - f_j \\
\mathcal{R}_2 &= \rho c_p \frac{\partial \hat{T}}{\partial t} - \frac{\partial}{\partial x_i}\left(k_{ij}\frac{\partial \hat{T}}{\partial x_j}\right) - Q - \beta_{ij}\hat{\sigma}_{ij}\frac{\partial \hat{\varepsilon}_{ij}}{\partial t} \\
\mathcal{R}_3 &= \hat{\sigma}_{ij} - C_{ijkl}\hat{\varepsilon}_{kl} + \beta_{ij}(\hat{T} - T_0)
\end{align}

\begin{equation}
\mathcal{L}_{PDE} = \frac{1}{N_{PDE}} \sum_{i=1}^{N_{PDE}} \left[ \mathcal{R}_1^2 + \mathcal{R}_2^2 + \mathcal{R}_3^2 \right]
\end{equation}

\begin{equation}
\mathcal{L}_{BC} = \frac{1}{N_{BC}} \sum_{i=1}^{N_{BC}} \left[ \|\hat{u} - u_{BC}\|^2 + \|\hat{T} - T_{BC}\|^2 + \|\hat{\sigma} \cdot n - \sigma_{BC}\|^2 \right]
\end{equation}

\begin{equation}
\mathcal{L}_{IC} = \frac{1}{N_{IC}} \sum_{i=1}^{N_{IC}} \left[ \|\hat{u}(t=0) - u_0\|^2 + \|\frac{\partial \hat{u}}{\partial t}(t=0) - \dot{u}_0\|^2 + \|\hat{T}(t=0) - T_0\|^2 \right]
\end{equation}

\subsection{Модуль предобработки}

\begin{itemize}
    \item Нормализация координат: $[-1, 1]$
    \item Интерполяция параметров среды
    \item Генерация коллокационных точек (LHS или адаптивно)
    \item Подготовка граничных условий
\end{itemize}

\subsection{Модуль обучения}

Оптимизация: Adam $\to$ L-BFGS

\begin{equation}
\theta_{k+1} = \theta_k - \alpha_k \nabla_\theta \mathcal{L}_{total}(\theta_k)
\end{equation}

Градиенты вычисляются через автоматическое дифференцирование.  

Адаптивное взвешивание:

\begin{equation}
\lambda_k(t) = \lambda_k^0 \cdot \exp\left(-\alpha_k \cdot \frac{\mathcal{L}_k(t)}{\mathcal{L}_{PDE}(t)}\right)
\end{equation}

\section{Особенности реализации для геологических сред}

\subsection{Учет неоднородности}

\begin{enumerate}
    \item Параметризация свойств:
    \begin{equation}
    \rho(x, y, z), \quad C_{ijkl}(x, y, z), \quad k_{ij}(x, y, z)
    \end{equation}
    \item Многослойная архитектура: отдельные подсети для различных геологических слоев
    \item Условные PINN: параметры среды как дополнительные входы
\end{enumerate}

\subsection{Обработка граничных условий}

Жесткие граничные условия:
\begin{equation}
\hat{u}_{hard} = g(x) + h(x) \cdot \hat{u}_{net}(x, t; \theta)
\end{equation}
где $g(x)$ удовлетворяет граничным условиям, $h(x) = 0$ на границе.  

Мягкие граничные условия включаются в функцию потерь.

\subsection{Масштабирование времени}

\begin{itemize}
    \item Разделение временного интервала на подынтервалы
    \item Обучение последовательности сетей
    \item Рекуррентные архитектуры (LSTM, GRU)
\end{itemize}

